\documentclass[x11names,dvipsnames,table]{beamer}
\usepackage{beamerthemesplit}
\usepackage{wrapfig}
\usetheme{SPbGU}
\usepackage{pdfpages}
\usepackage{amsmath}
\usepackage{cmap} 
\usepackage[T2A]{fontenc} 
\usepackage[utf8]{inputenc}
\usepackage[english,russian]{babel}
\usepackage{indentfirst}
\usepackage{amsmath}
\usepackage{tikz}
\usepackage{multirow}
\usepackage[noend]{algpseudocode}
\usepackage{algorithm}
\usepackage{algorithmicx}
\usetikzlibrary{shapes,arrows}
\usetikzlibrary{automata, arrows.meta, positioning}
\usetikzlibrary{decorations.pathreplacing, calligraphy}
\usepackage{fancyvrb}
\usepackage{colortbl}

\newtheorem{rutheorem}{Теорема}
\newtheorem{ruproof}{Доказательство}
\newtheorem{rudefinition}{Определение}
\newtheorem{rulemma}{Лемма}

\beamertemplatenavigationsymbolsempty


\title[]{Field-sensitive Points-to Analysis with CFL-Reachability}
% \subtitle[]{}
% То, что в квадратных скобках, отображается в левом нижнем углу. 
\institute[SPbU]{
JetBrains Research, Programming Languages and Tools Lab  \\
Saint Petersburg State University }

% То, что в квадратных скобках, отображается в левом нижнем углу.
\author[Vladimir Kutuev]{Vladimir Kutuev
}

\date{17.12.2021}

\begin{document}
{
% Лого университета или организации, отображается в шапке титульного листа
\begin{frame}
    \begin{tabular}{p{2.0cm} p{7.5cm} p{1cm}}
       \begin{center}
          \includegraphics[height=1.5cm]{pictures/jetbrainsResearch.pdf}
        \end{center}
        &
        \begin{center}
          \includegraphics[height=1.5cm]{pictures/YC_logo.pdf}
        \end{center}
        &
        \begin{center}
          \includegraphics[height=1.5cm]{pictures/SPbGU_Logo.png}
        \end{center}
      \end{tabular}
    \titlepage
\end{frame}
}

\begin{frame}[fragile]
  \transwipe[direction=90]
  \frametitle{CFL-Reachability}
  \begin{itemize}
    \item Context-free grammar --- $G = ( \Sigma, N, P, S )$
    \begin{itemize}
        \item $\Sigma$ --- set of terminal symbols
        \item $N$ --- set of non-terminal symbols
        \item $P$ --- set of production rules $\{A \to \alpha | A \in N, \alpha \in (\Sigma \cup N)^*\}$
        \item $S \in N$ --- start symbol
    \end{itemize}
    \item Graph --- $\mathcal{G}=(V,E,l)$
    \begin{itemize}
        \item $V$~--- set of vertices
        \item $E \subseteq V \times V$ --- set of edges 
        \item $l:E \to L$
    \end{itemize}
    \item $\{(v_1, v_n): \exists p=(e_1, e_2,..., e_n)) \in E^*, src(e_1)=v_1, dst(e_n)=v_n, l(e_1) \cdot...\cdot l(e_n) \in L(G)\}$
  \end{itemize}
\end{frame}
            
\begin{frame}[fragile]
\transwipe[direction=90]
\frametitle{Field-sensitive Points-to Analysis}

\begin{itemize}
    \item Graph \\
    \begin{tabular}{lll}
        Relation & Statement & Graph edge \\
        $alloc$\ $\subseteq$ Vars $\times$ Objects & x = new Obj(); & $x \xrightarrow{alloc} h$ \\
        $assign$\ $\subseteq$ Vars $\times$ Vars & x = y; & $x \xrightarrow{assign} y$ \\
        $load_f$\ $\subseteq$ Vars $\times$ Vars & x = y.f; & $x \xrightarrow{load_f} y$ \\
        $store_f$\ $\subseteq$ Vars $\times$ Vars & x.f = y; & $x \xrightarrow{store_f} y$ \\
    \end{tabular}

    \item Context-free grammar defining analysis \\
    \begin{tabular}{ll}
        $PointsTo\ \to\ (assign\ |\ load_f\ Alias\ store_f)^*$ $alloc$ & \\
        $Alias$ $\to$ $PointsTo$ \ $FlowsTo$ & $\forall$ f $\in$ $Fields$ \\
        $FlowsTo\ \to\ \overline{alloc}\ (\overline{assign}\ |\ \overline{store_f}\ Alias\ \overline{load_f})^*$ & \\
    \end{tabular}

\end{itemize}

\end{frame}


\begin{frame}[fragile]
\transwipe[direction=90]
\frametitle{Example Graph}

\begin{tabular}{cc}
    
    \begin{minipage}{.3\textwidth}
    \begin{Verbatim}[commandchars=\\\{\}]
v1 = new Obj(); // h1
v2 = new Obj(); // h2
v4 = new Obj(); // h3
v6 = new Obj(); // h4
v5 = v4;
v5 = v6;
v1.h = v1;
v2.g = v1;
v4.f = v2;
v7 = v5.f;
v9 = v6.f;
v3 = v2.g;
v8 = v7.g;
v10 = v9.g;
v10 = v10.h;
    \end{Verbatim}

\end{minipage}

    &
    \scalebox{0.8}{
\begin{minipage}{.4\textwidth}
    \begin{tikzpicture}[node distance=2.0cm] 
        \node[] (1) {$v5$};
        \node[] (2) [right of=1] {$v6$}; 
        \node[] (3) [left of=1] {$v4$};
        \node[] (4) [left of=3] {$h3$};
        \node[] (5) [right of=2] {$h4$};
        \node[] (6) [below of=3] {$v2$};
        \node[] (7) [left of=6] {$h2$};
        \node[] (8) [below of=2] {$v7$};
        \node[] (9) [below of=5] {$v9$};
        \node[] (10) [below of=6] {$v1$};
        \node[] (11) [left of=10] {$h1$};
        \node[] (12) [right of=10] {$v3$};
        \node[] (13) [below of=8] {$v8$};
        \node[] (14) [below of=9] {$v10$};
        \draw[->] (2) -- node[midway, above, sloped] {$alloc$} (5); 
        \draw[->] (3) -- node[midway, above, sloped] {$alloc$} (4); 
        \draw[->] (6) -- node[midway, above, sloped] {$alloc$} (7); 
        \draw[->] (10) -- node[midway, above, sloped] {$alloc$} (11); 
        \draw[->] (1) -- node[midway, above, sloped] {$assign$} (2); 
        \draw[->] (1) -- node[midway, above, sloped] {$assign$} (3);
        \draw[->] (14) to [out=-120, in=-60, looseness=3] node[midway, below] {$load_h$} (14); 
        \draw[->] (14) -- node[midway, right] {$load_g$} (9);
        \draw[->] (13) -- node[midway, right] {$load_g$} (8);
        \draw[->] (12) -- node[midway, right] {$load_g$} (6);
        \draw[->] (8) -- node[midway, right] {$load_f$} (1);
        \draw[->] (9) -- node[midway, right] {$load_f$} (2);
        \draw[->] (3) -- node[midway, right] {$store_f$} (6);
        \draw[->] (6) -- node[midway, left] {$store_g$} (10);
        \draw[->] (10) to [out=-60, in=-120, looseness=3] node[midway, below] {$store_h$} (10); 
    \end{tikzpicture} 
    \end{minipage}
    }

\end{tabular}

\end{frame}

\begin{frame}[fragile]
\transwipe[direction=90]
\frametitle{Example Graph with PointsTo Edges}

\begin{tabular}{cc}
    
    \begin{minipage}{.3\textwidth}
    \begin{Verbatim}[commandchars=\\\{\}]
v1 = new Obj(); // h1
v2 = new Obj(); // h2
v4 = new Obj(); // h3
v6 = new Obj(); // h4
v5 = v4;
v5 = v6;
v1.h = v1;
v2.g = v1;
v4.f = v2;
v7 = v5.f;
v9 = v6.f;
v3 = v2.g;
v8 = v7.g;
v10 = v9.g;
v10 = v10.h;
    \end{Verbatim}

\end{minipage}

    &
    \scalebox{0.8}{
\begin{minipage}{.4\textwidth}
    \begin{tikzpicture}[node distance=2.0cm] 
        \node[] (1) {$v5$};
        \node[] (2) [right of=1] {$v6$}; 
        \node[] (3) [left of=1] {$v4$};
        \node[] (4) [left of=3] {$h3$};
        \node[] (5) [right of=2] {$h4$};
        \node[] (6) [below of=3] {$v2$};
        \node[] (7) [left of=6] {$h2$};
        \node[] (8) [below of=2] {$v7$};
        \node[] (9) [below of=5] {$v9$};
        \node[] (10) [below of=6] {$v1$};
        \node[] (11) [left of=10] {$h1$};
        \node[] (12) [right of=10] {$v3$};
        \node[] (13) [below of=8] {$v8$};
        \node[] (14) [below of=9] {$v10$};
        \draw[->] (2) -- node[midway, above, sloped] {$alloc$} (5); 
        \draw[->] (3) -- node[midway, above, sloped] {$alloc$} (4); 
        \draw[->] (6) -- node[midway, above, sloped] {$alloc$} (7); 
        \draw[->] (10) -- node[midway, above, sloped] {$alloc$} (11); 
        \draw[->] (1) -- node[midway, above, sloped] {$assign$} (2); 
        \draw[->] (1) -- node[midway, above, sloped] {$assign$} (3);
        \draw[->] (14) to [out=-120, in=-60, looseness=3] node[midway, below] {$load_h$} (14); 
        \draw[->] (14) -- node[midway, right] {$load_g$} (9);
        \draw[->] (13) -- node[midway, right] {$load_g$} (8);
        \draw[->] (12) -- node[midway, right] {$load_g$} (6);
        \draw[->] (8) -- node[midway, right] {$load_f$} (1);
        \draw[->] (9) -- node[midway, right] {$load_f$} (2);
        \draw[->] (3) -- node[midway, right] {$store_f$} (6);
        \draw[->] (6) -- node[midway, left] {$store_g$} (10);
        \draw[->] (10) to [out=-60, in=-120, looseness=3] node[midway, below] {$store_h$} (10); 
        
        \draw[dashed,->,red] (3) to [out=135, in=45, looseness=1] (4);
        \draw[dashed,->,red] (1) to [out=120, in=60, looseness=1] (4);
        \draw[dashed,->,red] (2) to [out=45, in=135, looseness=1] (5);
        \draw[dashed,->,red] (1) to [out=60, in=120, looseness=1] (5);
        \draw[dashed,->,red] (6) to [out=135, in=45, looseness=1] (7);
        \draw[dashed,->,red] (8) to [out=135, in=45, looseness=1] (7);
        \draw[dashed,->,red] (10) to [out=135, in=45, looseness=1] (11);
        \draw[dashed,->,red] (12) to [out=-135, in=-45, looseness=1] (11);
        \draw[dashed,->,red] (13) to [out=-120, in=-60, looseness=1] (11);
    \end{tikzpicture} 
    \end{minipage}
    }

\end{tabular}

\end{frame}

\begin{frame}
  \transwipe[direction=90]
  \frametitle{Algorithms Based on Linear Algebra Operations}
  \begin{itemize}
    \item Matrix multiplication based algorithm
    \begin{itemize}
        \item Grammar should be converted to a Weak Chomsky Normal Form which significantly increases its size
        \item The performance of the algorithm depends on the grammar size
    \end{itemize}
    \item Kronecker Product based algorithm
    \begin{itemize}
        \item Grammar should be converted to a recursive state machine (RSM)
        \item For the intersection of RSM and a directed labeled graph, they are represented as a composition of adjacency boolean matrices for each label
    \end{itemize}
  \end{itemize}

\end{frame}

\begin{frame}
  \transwipe[direction=90]
  \frametitle{RSM for Field-sensitive Points-to Analysis}
  \begin{tabular}{|c|c|c|}
  \hline
  PointsTo & FlowsTo & Alias \\
    % Программа & Граф выражений \\
    
    \begin{minipage}{.36\textwidth}
    \scalebox{0.6}{
    \begin{tikzpicture}[node distance=1.9cm, auto] 
        \node (q0) [state, initial, initial text = {}] {$q_0$};
        \node (q1) [state, accepting, below = of q0] {$q_1$};
        
        \node (q4) [right = of q0] {$...$};
        
        \node (q2) [state, above = of q4] {};
        \node (q3) [state, right = of q2] {};
        
        \node (q5) [state, below = of q4] {};
        \node (q6) [state, right = of q5] {};
        
        \path [-stealth, thick]
            (q0) edge [loop above]  node[above left] {$assign$}()
            (q0) edge node[left] {$alloc$} (q1)
            
            (q0) edge[bend left] node[below right] {$load_{f_1}$} (q2)
            (q2) edge node[above] {$Alias$} (q3)
            (q3) edge node[below right] {$store_{f_1}$} (q0)
            
            (q0) edge[bend right] node[above right] {$load_{f_n}$} (q5)
            (q5) edge node[above] {$Alias$} (q6)
            (q6) edge node[above right] {$store_{f_n}$} (q0)
            ;
    \end{tikzpicture}
    }
    \end{minipage}

    &
\begin{minipage}{.36\textwidth}
\scalebox{0.6}{
    \begin{tikzpicture}[node distance=1.9cm] 
        \node (q0) [state, initial, initial text = {}] {$q_2$};
        \node (q1) [state, accepting, below = of q0] {$q_3$};
        
        \node (q4) [right = of q1] {$...$};
        
        \node (q2) [state, above = of q4] {};
        \node (q3) [state, right = of q2] {};
        
        \node (q5) [state, below = of q4] {};
        \node (q6) [state, right = of q5] {};
        
        \path [-stealth, thick]
            (q0) edge node[left] {$\overline{alloc}$} (q1)
            (q1) edge [loop below] node[below left] {$\overline{assign}$}()
            
            (q1) edge[bend left] node[below right] {$\overline{store_{f_1}}$} (q2)
            (q2) edge node[above] {$Alias$} (q3)
            (q3) edge node[below right] {$\overline{load_{f_1}}$} (q1)
            
            (q1) edge[bend right] node[above right] {$\overline{store_{f_n}}$} (q5)
            (q5) edge node[above] {$Alias$} (q6)
            (q6) edge node[above right] {$\overline{load_{f_n}}$} (q1);
    \end{tikzpicture} 
    }
    \end{minipage}
    &
\begin{minipage}{.17\textwidth}
\scalebox{0.7}{
    \begin{tikzpicture}[node distance=1.0cm] 
        \node (q0) [state, initial, initial text = {}] {$q_4$};
        \node (q1) [state, below = of q0] {$q_5$};
        \node (q2) [state, accepting, below = of q1] {$q_6$};
        \path [-stealth, thick]
            (q0) edge node[right] {$PointsTo$} (q1)
            (q1) edge node[right] {$FlowsTo$} (q2)
            ;
    \end{tikzpicture} 
    }
    \end{minipage}
    \\
    & & \\
    \hline

\end{tabular}

\end{frame}

\begin{frame}
  \transwipe[direction=90]
  \frametitle{One Terminal, One Non-terminal Representation}
  \begin{itemize}
    \item The graph and RSM are represented as adjacency matrices, whose elements are integers
    \item The high-order bits encode the non-terminal, the low-order bits encode the terminal \\
  \end{itemize}
  \centering {
    \scalebox{0.7}{
      \begin{tikzpicture}
        [%%%%%%%%%%%%%%%%%%%%%%%%%%%%%%
            box/.style={rectangle,draw=black, ultra thick, minimum size=1cm},
        ]%%%%%%%%%%%%%%%%%%%%%%%%%%%%%%
        \foreach \x in {0,1,2,3,4,5,6,7,8}
            \node[box] at (\x,0){};
        \draw[decorate,decoration={brace, amplitude=0.2cm,mirror},thick] (-.5,-.7) -- node[below]{Non-tetminal} (1.49,-.7);
        \draw[decorate,decoration={brace, amplitude=0.2cm, mirror},thick] (1.51,-.7) -- node[below]{Terminal} (8.5,-.7);
      \end{tikzpicture}
    }
  }
  \begin{itemize}
    \item Operation for Kronecker product
    \begin{itemize}
        \item $times(x, y) = (x_{nonterm} = y_{nonterm} \neq 0 \ or\ x_{term} = y_{term} \neq 0)$
    \end{itemize}
  \end{itemize}
\end{frame}

\begin{frame}
  \transwipe[direction=90]
  \frametitle{Adapting Graph to this Representation}
  \begin{tabular}{ccc}
  Edges & New variable & Result \\
 
  \begin{tikzpicture}[node distance=2.0cm] 
        \node[] (1) {$x$};
        \node[] (2) [right of=1] {$y$};
        
        \draw[->] (1) to [out=45, in=135, looseness=1] node[midway, above] {$...$} (2);
        \draw[->] (1) to node[midway, below, sloped] {$assign$} (2); 
  \end{tikzpicture} 
  &
  z = y; x = z;
  & 
  \begin{tikzpicture}[node distance=2.0cm] 
        \node[] (1) {$x$};
        \node[] (2) [right of=1] {$z$};
        \node[] (3) [right of=2] {$y$};
        
        \draw[->] (1) -- node[midway, above, sloped] {$assign$} (2); 
        \draw[->] (2) -- node[midway, above, sloped] {$assign$} (3); 
        \draw[->] (1) to [out=45, in=135, looseness=1] node[midway, above] {$...$} (3);
  \end{tikzpicture} 
  \\
  \begin{tikzpicture}[node distance=2.0cm] 
        \node[] (1) {$x$};
        \node[] (2) [right of=1] {$y$};
        
        \draw[->] (1) to [out=45, in=135, looseness=1] node[midway, above] {$...$} (2);
        \draw[->] (1) to node[midway, below, sloped] {$store_f$} (2); 
  \end{tikzpicture} 
  &
  z = y; x.f = z;
  & 
  \begin{tikzpicture}[node distance=2.0cm] 
        \node[] (1) {$x$};
        \node[] (2) [right of=1] {$z$};
        \node[] (3) [right of=2] {$y$};
        
        \draw[->] (1) -- node[midway, above, sloped] {$store_f$} (2); 
        \draw[->] (2) -- node[midway, above, sloped] {$assign$} (3); 
        \draw[->] (1) to [out=45, in=135, looseness=1] node[midway, above] {$...$} (3);
  \end{tikzpicture} 
  \\
  \begin{tikzpicture}[node distance=2.0cm] 
        \node[] (1) {$x$};
        \node[] (2) [right of=1] {$y$};
        
        \draw[->] (1) to [out=45, in=135, looseness=1] node[midway, above] {$...$} (2);
        \draw[->] (1) to node[midway, below, sloped] {$load_f$} (2); 
  \end{tikzpicture} 
  &
  z = y.f; x = z;
  & 
  \begin{tikzpicture}[node distance=2.0cm] 
        \node[] (1) {$x$};
        \node[] (2) [right of=1] {$z$};
        \node[] (3) [right of=2] {$y$};
        
        \draw[->] (1) -- node[midway, above, sloped] {$assign$} (2); 
        \draw[->] (2) -- node[midway, above, sloped] {$load_f$} (3); 
        \draw[->] (1) to [out=45, in=135, looseness=1] node[midway, above] {$...$} (3);
  \end{tikzpicture} 
  \end{tabular}
\end{frame}


\begin{frame}
  \transwipe[direction=90]
  \frametitle{Results}
  \begin{itemize}
    \item The algorithm based on the Kronecker product has been optimized for the field-sensitive points-to analysis problem 
    \item The developed algorithm was implemented within the CFPQ\_PyAlgo\footnote{\url{https://github.com/JetBrains-Research/CFPQ_PyAlgo}} framework
    \item The correctness of the algorithm has been checked on the Dataset\footnote{\url{https://bitbucket.org/jensdietrich/gigascale-pointsto-oopsla2015}}
  \end{itemize}

  \textbf{Future Work}:
  \begin{itemize}
    \item Extend Benchmark Graphalytics\footnote{\url{https://graphalytics.org/}}
    \begin{itemize}
        \item Investigate the performance of the developed algorithm
        \item Compare performance with other CFL-Reachability algorithms 
    \end{itemize}
  \end{itemize}
\end{frame}


\end{document}
